\section{问题三模型、求解与结果}

\subsection{滚动信息时间线}

令更新时刻$k\in\{0,6,12,18\}$ h，$m(k)=6k+1$。在$k$时刻仅优化尚未执行的时段$\mathcal T_k=\{m(k),\ldots,144\}$，并只执行到下一次更新。信息集包含当前真实SOC、已实现的负荷与光伏、已发布预报、固定电价和0:00原计划，不包含未来实际值或之后才发布的预报。

对发布时间$a$、当前时刻$\tau$和目标时刻$v$，预测提前量$h=v-a$，预报年龄$A=\tau-a$，距目标时间$H=v-\tau$，满足$h=A+H$。使用旧预报时不能把$H$误当作该预报原始提前量。附件3的小时节点按线性函数在10 min区间积分，形成与公共模型一致的kWh输入。

\subsection{提前量校准、联合场景与可信度指标}

主方案以过去28个完整日估计分发布时间、6 h提前量箱的偏差与尺度，并用3日收缩稳定小样本估计。候选还包括统一尺度和指数尺度
\begin{equation}
 \sigma_g(h)=\sigma_{g,0}\exp(\beta h),\qquad \beta\ge0,
\end{equation}
以及归一化展示指标
\begin{equation}
 R_g(h)=\frac{1}{1+\{\sigma_g(h)/s_{\mathrm{ref}}\}^2}.
\end{equation}
$R$只反映条件误差尺度并可反解场景离散度，不是预测正确概率，也不改变预测中心或场景概率。场景按同一历史日联合抽取负荷与光伏标准化残差块，每次64个抽样；各更新时刻根据当前已发布预报和真实SOC重建剩余时域场景。

\subsection{调整结算与CVaR随机MPC}

令$q_t$为0:00原计划，$r_{t|k}$为更新后的正常购电量。主结算为
\begin{equation}
 g_t(q,r)=p_tq+1.5p_t\pos{r-q}-0.5p_t\pos{q-r}
 =p_tr+0.5p_t|r-q|.
\end{equation}
它表示相对原计划的上调和下调均承担50\%偏差成本。定义剩余时域场景紧急费用$C_k^\omega=\sum_{t\in\mathcal T_k}5p_te_{t|k}^\omega$，风险增强目标为
\begin{equation}
 \min\ \sum_{t\in\mathcal T_k}g_t(q_t,r_{t|k})+\E[C_k]
 +\lambda\operatorname{CVaR}_{\alpha}(C_k)+\gamma_k b_k^{\mathrm{end}},
\end{equation}
其中
\begin{equation}
 \operatorname{CVaR}_{\alpha}(C_k)=\min_{\zeta}\left\{\zeta+
 \frac{1}{1-\alpha}\sum_\omega\pi_\omega\pos{C_k^\omega-\zeta}\right\}.
\end{equation}
CVaR表示剩余紧急费用的上尾风险\cite{rockafellar2000}；风险惩罚和日末软目标只参与决策，不计入实际账单。每天在四个更新点重算电池计划，区间内固定执行到下一节点；全年末统一回到6000 kWh。

\subsection{消融、预测检验与经济结果}

表\ref{tab:q3-ablation}按相同334天、相同跨日衔接给出主要消融。B0仍保留每日四次求解，但只让0:00预报成为新信息；B1使用官方点预测；B2、B3依次加入联合场景和提前量条件分箱尺度；B4在B3上增加CVaR。由此避免把控制机会与新信息价值混为一谈。

\begin{table}[H]
\centering
\caption{问题三B0--B4主要对照}
\label{tab:q3-ablation}
\begin{tabularx}{\textwidth}{Xrrr}
\toprule
方案 & 总费/万元 & 紧急费/万元 & 日紧急费CVaR95/元\\
\midrule
B0 仅0:00新预报 & 1703.98 & 126.22 & 17898.31\\
B1 官方点预测 & 2401.59 & 1196.73 & 80723.79\\
B2 统一尺度联合场景 & 1662.76 & 91.56 & 16854.47\\
B3 条件分箱尺度（主方案） & 1662.77 & 91.25 & 16718.04\\
B4 CVaR90，$\lambda=0.1$ & 1671.75 & 65.99 & 13088.86\\
\bottomrule
\end{tabularx}
\end{table}

B3实际总费为16627713.37元，其中紧急购电费912512.24元、紧急购电量238.47 MWh。B2费用比B3低92.03元，表明条件尺度在本年度并未稳定降低平均费用；但B3的紧急费和经验CVaR95略低，故不能把该模块宣称为全面优于统一尺度。B4将日紧急费CVaR95降至13088.86元、紧急费降至659931.40元，同时使总费增加89781.36元，体现平均费用与尾部风险之间的明确权衡。

光伏正值样本的合并CRPS为184.54 kW，90\%经验区间覆盖率为80.67\%。覆盖率低于标称值，说明有限日块及非负截断后的区间仍偏窄；$R$与$\sigma$的一一映射不能提供概率安全保证。

\begin{figure}[H]
  \centering
  \includegraphics[width=0.92\textwidth]{q3_ablation.png}
  \caption{问题三B0--B4费用与风险对照}
  \label{fig:q3-ablation}
\end{figure}

\subsection{预报更新时间价值}

全四次新预报相对仅0:00新预报累计节省412063.77元。已有预报时刻的8种组合中，事后最低费用组合为0:00、6:00和12:00，共节省415281.22元；加入18:00预报的边际节省为$-3217.45$元。配对7日移动块Bootstrap给出的18:00边际日节省95\%探索性区间为$[-21.41,0.78]$元，不能据此断言该时刻必然有正价值。附件未提供3:00、9:00、15:00和21:00的新预报，重算旧预报不能冒充额外气象信息。

\begin{figure}[H]
  \centering
  \includegraphics[width=0.88\textwidth]{q3_information_value.png}
  \caption{问题三已有预报时刻的增量信息价值}
  \label{fig:q3-information}
\end{figure}

主方案共完成1336次外层滚动、334天执行回放；全部对照共含21组334日实验、28056次正式滚动求解。物理约束、信息时序、跨日连续、结算重算和概率尺度均经独立审计，最大残差为$7.55\times10^{-7}$，20项自动测试通过。

\section{问题四模型与结果占位}

\subsection{价格信息与共同结算口径}

问题四沿用式\eqref{eq:balance}--\eqref{eq:soc}。主口径按交付时段的实际波动价结算；预测价格只服务决策，不能替代实际账单。Q4-2每日0:00依据历史价格和历史负荷--光伏信息制定全天计划；Q4-3在0:00、6:00、12:00和18:00使用届时已发布的光伏预报、已实现价格和真实SOC。若把当天未来实际价格提前输入，只能作为完美信息下界。

令净负荷$N_{i,t}=L_{i,t}-G_{i,t}$。基于历史滚动样本构造价格与净负荷的成对场景$\{p_{i,t}^{\omega},N_{i,t}^{\omega}\}$，并以简单日历分组价格基线和已实现残差修正作为冷启动友好的预测。实际账单分别为
\begin{align}
 F_{4,2}^{\mathrm{real}}&=\sum_t p_{i,t}^{\mathrm{real}}(q_{i,t}+5e_{i,t}^{\mathrm{real}}),\\
 F_{4,3}^{\mathrm{real}}&=\sum_t p_{i,t}^{\mathrm{real}}\{b(q_{i,t},r_{i,t}^{\mathrm{exec}})+5e_{i,t}^{\mathrm{real}}\},
\end{align}
其中$b(q,r)=r+0.5|r-q|$。

\subsection{Q4-2：价量依赖收缩}

小样本下直接使用经验成对分布可能过度估计价量依赖。以历史验证确定全局收缩权重$\rho\in[0,1]$：
\begin{equation}
 \widehat{\mathbb P}^{\rho}_t=\rho\widehat{\mathbb P}^{\mathrm{pair}}_t+
 (1-\rho)(\widehat{\mathbb P}^{p}_t\otimes\widehat{\mathbb P}^{N}_t).
\end{equation}
相应的价格加权缺电费用曲线为
\begin{equation}
 H_t^\rho(A)=\sum_\omega a_{t,\omega}^\rho(N_t^\omega-A)^+,qquad
 a_{t,\omega}^\rho=\pi_\omega\{\rho p_t^\omega+(1-\rho)\bar p_t\}.
\end{equation}
$\rho=1$恢复成对模型，$\rho=0$保留边际但假定价量独立。$\rho$不是相关系数、预测正确概率或当天优化变量；实际执行仍按真实价格收费。Q4-2在公共约束下最小化正常购电期望费用、$5H_t^\rho$对应的紧急费用及终端价值，并保留$\rho=0,1$端点对照。

\subsection{Q4-3：前瞻下一次更新的条件调整}

在当前更新点$k$，只展开下一节点$\kappa$，以届时可观测的“新旧光伏预报修订”和节点前已实现价格偏差构造$K\in\{1,2,3\}$个信息组。分叉前的正常购电、充放电与库存对所有场景共享；各分支从同一桥接库存$S_b$出发，分叉后动作只允许依据信息组变化：
\begin{equation}
 S_b^m=S_b^0,qquad m=1,\ldots,K.
\end{equation}
目标联合选择0:00原承诺$q$、当前共同动作和下一节点后的条件调整方案，使未来继续运行价值反过来影响当前留存电量。组号只由下一节点真实可见的信号决定，不能使用其后的实际负荷、实际价格或最优费用。该结构受分桶非预见性思想启发\cite{li2016}；当$K=1$且场景、终端和结算一致时，应恢复不前瞻分组的旧版基准。

\subsection{待完成结果与验收}

\todo{完成Q4-2的$\rho=0$、$\rho=1$和历史验证所选$\rho$对照，以及Q4-3的$K=1$退化核验和$K=2,3$条件前瞻对照；填写波动电价下的费用、紧急购电、套利收益、典型日、完美信息下界和敏感性结果。}

在数值结果完成前，仅保留如下验收标准：收缩权重和为$\bar p_t$且单位为元/kWh；全部分支共享桥接库存；分组映射不泄漏未来；各方案用相同日期、控制机会和连续库存回放真实账单；不得以优化器内部预测目标下降代替真实样本外收益。
